\begin{tikzpicture}[line cap=round, line join=round, font=\small, scale=0.65]
	\coordinate (L1) at (-3,3);
	\coordinate (L2) at (-3,2.5);
	\coordinate (L3) at (-6,2.5);
	\coordinate (L4) at (-6,2);
	\coordinate (L5) at (0,2);
	\coordinate (L6) at (0,3);

	\node at (-3,5) {$\Omega$};

	\path[fill=lightgray] (L6)--(L1)--(L2)--(L3)--(L4)--(L5)--cycle;
	\path[fill=darkgray, pattern=north east lines, pattern color=darkgray]
	(L6)--(L1)--(L2)--(L3)--(L4)--(L5)--cycle;

	\draw[black, thin] (L3)--(L4)--(L5)--(L6);
	\draw[black, ultra thick] (L3)--(L2)--(L1)--(L6);

	\coordinate (A) at (-1,7.05);
	\coordinate (B) at (-1,3.05);
	\coordinate (C) at (-5,7.05);
	\coordinate (D) at (-5,3.05);
	\coordinate (E) at (-1,3.05);

	\begin{scope}[on background layer]
		\path[fill=lightgray, opacity=0.3] (A)--(B)--(D)--(C)--cycle;
	\end{scope}

	\draw[draw=red, ultra thick] (A) -- (B)
	node[midway, right, text=red] {$\Gamma_D$};

	\draw[draw=blue, ultra thick] (A) -- (C)
	node[midway, above, text=blue] {$\Gamma_N$};

	\draw[draw=blue, ultra thick] (C) -- (D)
	node[midway, left, text=blue] {$\Gamma_N$};

	% Pfeil zuerst
	\draw[orange, thick, -{Latex[length=1mm,width=1mm,fill=orange]}]
	(-2,3.05) -- (-2,2.25)
	node[midway, right, text=orange] {$\mathrm{n}$};

	% Danach Gamma_C, damit die Linie über dem Pfeilanfang liegt
	\draw[draw=teal, ultra thick] (D) -- (E)
	node[midway, above, text=teal] {$\Gamma_C$};

\end{tikzpicture}